\subsection{Floating-Point Exceptions, Conversions, and Reductions}

Each active FP lane applies the corresponding scalar rounding, NaN,
signed-zero, default-NaN, denormal, and exception rules.  Inactive lanes
generate no cause.  The instruction unions all active-lane causes before
testing the enables in the original \texttt{FSTATUS}.  If any generated cause
is enabled, one (:ref:FP.events.EXCEPTION.FLOATING_POINT_EXCEPTION:) occurs and neither a destination
effect nor newly accrued \texttt{FFLAGS} becomes visible.  Otherwise the union
is accrued once and all destination effects commit together.  FP-to-integer
conversion rounds toward zero and follows the scalar saturation and invalid
result rules.

Integer reductions return a zero-extended result in an \texttt{Rn}.  For an
empty predicate, ADD, OR, and XOR return zero; AND and unsigned MIN return the
all-ones value of the selected width; signed MIN returns the signed maximum;
signed MAX returns the signed minimum; and unsigned MAX returns zero.  FP
reductions return an \texttt{Fn}.  Empty FP ADD returns positive zero, empty FP
MIN returns positive infinity, and empty FP MAX returns negative infinity.
FP ADD visits active lanes in increasing logical-lane order and may not be
reassociated.
